\documentclass[11pt,a4paper]{article}

\usepackage{fontspec}
\usepackage{xeCJK}
\setmainfont{TeX Gyre Heros}
\setsansfont{TeX Gyre Heros}
\setmonofont{DejaVu Sans Mono}
\setCJKmainfont[AutoFakeBold=2,AutoFakeSlant=0.2]{Droid Sans Fallback}
\setCJKsansfont[AutoFakeBold=2,AutoFakeSlant=0.2]{Droid Sans Fallback}
\setCJKmonofont{Droid Sans Fallback}
\XeTeXlinebreaklocale "zh"
\XeTeXlinebreakskip = 0pt plus 1pt

\usepackage[margin=2.25cm,headheight=15pt]{geometry}
\usepackage{amsmath,amssymb,amsthm}
\usepackage{array,booktabs,longtable,tabularx}
\usepackage{enumitem}
\usepackage{xcolor}
\usepackage{graphicx}
\usepackage{fancyhdr}
\usepackage{setspace}
\usepackage{hyperref}
\makeatletter
\g@addto@macro\UrlBreaks{\do\_\do\-}
\makeatother

\definecolor{WanBlue}{HTML}{235789}
\definecolor{WanLightBlue}{HTML}{EAF2F8}
\definecolor{WanGreen}{HTML}{18794E}
\definecolor{WanLightGreen}{HTML}{EAF7F1}
\definecolor{WanOrange}{HTML}{A64B00}
\definecolor{WanLightOrange}{HTML}{FFF1E6}
\definecolor{WanGray}{HTML}{5B6573}
\definecolor{WanLightGray}{HTML}{F4F6F7}

\hypersetup{
  unicode=true,
  colorlinks=true,
  linkcolor=WanBlue,
  citecolor=WanGreen,
  urlcolor=WanBlue,
  pdftitle={CPS 训练模型在 CPS 与 ODE 评测间的性能差异},
  pdfauthor={Wan-Trainer experiment audit},
  pdfsubject={Flow-CPS, ODE sampling, DanceGRPO, VBVR-Pro EvalKit}
}

\pagestyle{fancy}
\fancyhf{}
\lhead{\small Wan-Trainer 技术报告}
\rhead{\small CPS、ODE 与 EvalKit 审计}
\cfoot{\thepage}

\setlength{\parindent}{2em}
\setlength{\parskip}{0.35em}
\setlength{\emergencystretch}{2em}
\setstretch{1.08}
\setlist{nosep,leftmargin=2em}
\renewcommand{\contentsname}{目录}
\renewcommand{\abstractname}{摘要}
\renewcommand{\refname}{参考文献}
\renewcommand{\tablename}{表}
\renewcommand{\figurename}{图}
\renewcommand{\proofname}{证明}

\newcommand{\code}[1]{\texttt{#1}}
\newcommand{\score}[1]{\ensuremath{\mathbf{#1}}}
\newcommand{\best}[1]{\textcolor{WanGreen}{\score{#1}}}
\newcommand{\warnbox}[1]{%
  \noindent\colorbox{WanLightOrange}{%
    \parbox{\dimexpr\textwidth-2\fboxsep\relax}{\textcolor{WanOrange}{#1}}%
  }%
}
\newcommand{\resultbox}[1]{%
  \noindent\colorbox{WanLightGreen}{%
    \parbox{\dimexpr\textwidth-2\fboxsep\relax}{\textcolor{WanGreen}{#1}}%
  }%
}
\newcommand{\infobox}[1]{%
  \noindent\colorbox{WanLightBlue}{%
    \parbox{\dimexpr\textwidth-2\fboxsep\relax}{\textcolor{WanBlue}{#1}}%
  }%
}

\newtheorem{proposition}{命题}
\newtheorem{observation}{观察}

\title{
  \vspace{-1.2cm}
  \textbf{\LARGE CPS 训练模型在 CPS 与 ODE 评测间的性能差异}\\[0.5em]
  \Large 论文、Wan-Trainer 实现与 VBVR-Pro EvalKit 更新审计
}
\author{Wan-Trainer 实验复盘}
\date{2026 年 7 月 24 日}

\begin{document}

\maketitle

\begin{abstract}
本文回答三个相互关联的问题：
第一，Wan-Trainer 中 Flow-CPS 的训练与生成实现是否忠实于
\emph{Coefficients-Preserving Sampling for Reinforcement Learning with Flow Matching}
的公式\cite{flowcps}；
第二，为什么同一个 CPS/DanceGRPO 检查点用 CPS 生成时的 VBVR-Pro 分数显著高于
ODE 生成；
第三，VBVR EvalKit 的 \code{main\_v2} 最近更新后，训练奖励与最终评测代码应如何同步。

审计结论是：核心 CPS 转移公式与论文一致，当前现象不是由一个明显的公式抄错造成的。
更主要的原因是\textbf{采样器依赖的强化学习目标}与\textbf{未严格配平的评测协议}：
历史检查点在 $T=30$、CFG $=1$ 的随机 Flow-CPS 轨迹上收集奖励并更新策略，
而常规 ODE 评测使用 $T=50$、CFG $=5$ 的 UniPC 或 FlowMatch Euler 轨迹。
即使把步数和 CFG 配平，用 $\eta=0$ 的 Flow-CPS 作为一阶 Euler 控制，
$\eta=0.7$ 仍在最新 EvalKit 上领先 $0.023754$，说明差异中存在稳定的
sampler-dependent performance。这一现象与 RL 后的 sampler co-adaptation
一致，但单个 post-RL 检查点尚不能把差值全部归因于 RL 优化。

在 500 个无错误样本、100 个任务、同一检查点 600 的最新 EvalKit
\code{main\_v2} 复评分中，CPS $\eta=0.7$ 的 Overall 为
\best{0.534130}，UniPC ODE 为 $0.504266$，FlowMatch Euler ODE 为
$0.504909$，CPS $\eta=0$ 为 $0.510376$。以任务为重采样单位的
100,000 次 task-paired bootstrap 显示，CPS $\eta=0.7$ 相对三种控制的
Overall 差值 95\% 区间均不跨零。不过，这仍是已有视频的 scorer-only
复评分，且跨采样器视频尚未证明使用了完全相同的初始 latent，因此不能写成最终因果证明。
\end{abstract}

\resultbox{\textbf{一句话结论：}
CPS 训练优化的是 ``模型 + CPS 转移核 + 训练时 CFG/时间网格'' 这一整套策略；
ODE 是另一套状态分布和转移核。当前 ODE 分数并非不可提高，但第一优先级应是
做同初始 latent、同时间网格、同 CFG 的受控比较，再决定是调 ODE 推理参数，
还是加入 ODE-aware 训练目标。}

\tableofcontents
\newpage

\section{问题、范围与证据等级}

\subsection{要回答的问题}

本文把原始现象拆成四个可检验的问题：

\begin{enumerate}
  \item \textbf{公式一致性：}本项目的 Flow-CPS 单步均值、噪声标准差和训练代理
  log-probability 是否与论文及官方实现一致？
  \item \textbf{边界一致性：}当 $\eta=0$ 时，Flow-CPS 是否应退化为 ODE；
  如果应当退化，需要满足哪些额外条件？
  \item \textbf{实验解释：}CPS $\eta=0.7$ 比 ODE 高，究竟是采样器特异优化，
  还是步数、CFG、求解器、随机种子或 scorer 版本混杂造成的？
  \item \textbf{工程闭环：}EvalKit 更新后，训练奖励、最终评测、历史报告和
  provenance 如何避免静默混用？
\end{enumerate}

\subsection{证据分层}

本文明确区分三类结论：

\begin{description}
  \item[代数结论] 可由公式直接推出，例如 $\eta=0$ 的 CPS 单步等价于显式 Euler。
  \item[实现结论] 来自源代码、配置和测试审计，例如项目实现了论文的系数分解，
  且 CPS 评测默认 $T=30$、CFG $=1$。
  \item[实验结论] 来自检查点 600 的 500 样本复评分和任务级 bootstrap；
  它支持 sampler-dependent gap，但不能替代严格的同 latent 生成实验，
  也不能单独证明该 gap 是 RL 才产生的。
\end{description}

\warnbox{\textbf{范围限制：}
本文重点分析 Wan2.2 TI2V-5B 的 strict In-Domain DanceGRPO 检查点 600。
这里的 ``Out-of-Domain'' 是 EvalKit 的任务分组名称；这些任务类型曾出现在本次
RL 训练数据中，只是 bench 样本 ID 与训练样本分离。因此不能将 OOD 分数描述为
零样本任务泛化。}

\section{论文中的 Flow-CPS：公式与正确解释}

\subsection{Rectified Flow 记号}

令 $x_0$ 表示干净数据 latent，$x_1\sim\mathcal{N}(0,I)$ 表示噪声，
采用 Rectified Flow 的线性插值\cite{rectifiedflow}
\begin{equation}
  x_t=(1-t)x_0+t x_1,\qquad t\in[0,1].
  \label{eq:rf-interpolation}
\end{equation}
由当前状态与模型预测速度 $v_\theta(x_t,t,c)$ 可定义两个预测端点
\begin{align}
  \widehat{x}_0 &= x_t-t\,v_\theta(x_t,t,c), \label{eq:x0hat}\\
  \widehat{x}_1 &= x_t+(1-t)\,v_\theta(x_t,t,c). \label{eq:x1hat}
\end{align}
只有当 $v_\theta$ 等于真实配对速度 $x_1-x_0$ 时，它们才分别重建
真实的 $x_0$ 与 $x_1$。
从较大噪声时刻 $t$ 走到较小噪声时刻 $s<t$ 时，标准一阶 ODE 更新为
\begin{equation}
  x_s^{\mathrm{Euler}}
  =x_t+(s-t)v_\theta(x_t,t,c).
  \label{eq:euler}
\end{equation}

\subsection{Flow-CPS 单步}

Flow-CPS 论文\cite{flowcps}以 $\eta\in[0,1]$ 控制随机强度，
并把下一时刻 $s$ 的总噪声系数
分解为预测噪声方向与新高斯噪声方向：
\begin{equation}
  x_s^{\mathrm{CPS}}
  =(1-s)\widehat{x}_0
  +s\cos\!\left(\frac{\pi\eta}{2}\right)\widehat{x}_1
  +s\sin\!\left(\frac{\pi\eta}{2}\right)z,
  \qquad z\sim\mathcal{N}(0,I).
  \label{eq:flow-cps}
\end{equation}
如果把条件均值与随机标准差分开写，则
\begin{align}
  \mu_\theta(x_t,t,s,\eta)
    &=(1-s)\widehat{x}_0
      +s\cos\!\left(\frac{\pi\eta}{2}\right)\widehat{x}_1,\\
  \sigma_{\mathrm{CPS}}(s,\eta)
    &=s\sin\!\left(\frac{\pi\eta}{2}\right).
\end{align}

这里的 ``coefficients-preserving'' 是一个\textbf{调度系数约束}。
算法按构造使新采样的 $z$ 独立于当前 $x_t$，因而也独立于作为
$x_t$ 确定函数的 $\widehat{x}_1$。若进一步把 $\widehat{x}_1$
视作单位高斯噪声方向，则组合后的名义总噪声标准差满足
\begin{equation}
  s\sqrt{
    \cos^2\!\left(\frac{\pi\eta}{2}\right)
    +\sin^2\!\left(\frac{\pi\eta}{2}\right)}
  =s.
\end{equation}
因此下一步的信号系数为 $1-s$、总噪声系数为 $s$，与原 scheduler 对齐。

\warnbox{\textbf{不要过度解读系数保持：}
$z$ 按构造是独立标准高斯，但真实 learned $\widehat{x}_1$ 不保证具有
零均值、单位协方差的 Gaussian 条件律，也不保证与
$\widehat{x}_0$ 形成理想的联合分解。故式 \eqref{eq:flow-cps}
锁定的是名义 coefficient schedule，而不是对有限步、误差模型给出
``CPS 与 probability-flow ODE 终点分布完全相同'' 的定理。}

最多只有在下列条件同时成立时，才可期待不同 $\eta$ 共享目标终点边缘分布：
(1) 速度/score 在全部 scheduler marginal 上精确；
(2) 预测噪声具有正确条件律，且新 $z$ 为独立标准高斯；
(3) 每个离散转移核本身满足
$p_s(x_s)=\int K_{\theta,\eta}(x_s\mid x_t)p_t(x_t)\,dx_t$，
而不只是二阶系数正确；
(4) conditioning、CFG、schedule 与时间参数化完全一致；
(5) 离散误差消失，或该核另有有限步 marginal-preservation 证明。
即使这些条件成立，结论也只是 distributional/marginal 一致，
不是同 seed 的路径一致。论文没有为一般的有限步 learned kernel 证明第 (3) 条。

\subsection{\texorpdfstring{$\eta=0$}{eta=0} 与 Euler 的严格关系}

\begin{proposition}
在实数算术下，当 $\eta=0$，且 CPS 与 ODE 使用同一个当前状态 $x_t$、
同一个 $v_\theta$、相同 conditioning/CFG 和相同 $(t,s)$ 时间网格时，
Flow-CPS 单步在代数上严格等于显式 Euler 单步；这不要求模型预测完美。
\end{proposition}

\begin{proof}
代入 $\cos 0=1$、$\sin 0=0$，并使用
式 \eqref{eq:x0hat}--\eqref{eq:x1hat}：
\begin{align}
  x_s^{\mathrm{CPS},\eta=0}
  &=(1-s)(x_t-t v_\theta)+s(x_t+(1-t)v_\theta)\\
  &=x_t+\bigl[-t(1-s)+s(1-t)\bigr]v_\theta\\
  &=x_t+(s-t)v_\theta
  =x_s^{\mathrm{Euler}}.
\end{align}
\end{proof}

浮点实现中，两条代数等价式可能采用不同运算顺序：
CPS 先形成 $\widehat{x}_0,\widehat{x}_1$ 再重组，而 direct Euler
直接计算 $x_t+(s-t)v_\theta$。因此工程验证应预先规定
\code{atol}/\code{rtol}，检查逐步误差及其累积；只有运算图、dtype、
autocast 与 rounding 也相同，才可要求 bitwise equality。

这个命题很有用，但它\textbf{不是} ``任意 ODE 评测都应与 CPS $\eta=0$ 相同''。
以下任一变化都会破坏代数等价所需的条件，或扩大数值误差：

\begin{itemize}
  \item ODE 使用 UniPC\cite{unipc} 等多步或 predictor-corrector 高阶算法；
  \item 两条路径使用不同 timestep/sigma 网格、flow shift 或末端处理；
  \item CFG 不同，导致每一步的有效速度场不同；
  \item VAE、condition 构造、dtype、随机初始 latent 或 seed-to-latent 映射不同；
  \item 一个路径通过 Diffusers pipeline，另一个走训练时内部 sampler，
  两者有不同的 scheduler 细节。
\end{itemize}

因此，CPS $\eta=0$ 是一个非常好的\textbf{同实现路径内的一阶 ODE 控制}，
但不能直接替代对另一个 pipeline 的容差内逐步数值一致性测试。

\subsection{Flow-CPS、Flow-CPWS 与 Flow-SDE 不能混为一谈}

论文第 4.4 节还定义了带 Wiener 缩放的 Flow-CPWS，并说明 Flow-SDE 是
Flow-CPWS 在 $\Delta t\to0$ 下的一阶 Taylor 近似，误差来自噪声系数展开。
论文给出的 noise-level error 形式为
\[
  \sqrt{
    \frac{(\sigma_t\Delta t)^2}{t}
    +\left(\frac{\sigma_t^2\Delta t}{2t}\right)^2
  },
\]
其中 $1/t$ 因子使 $t\to0$ 附近即使步数较多也可能恶化。
Theorem 1 的对象是
\[
  \text{Flow-SDE}\ \approx\ \text{Flow-CPWS},
\]
而不是证明
\[
  \text{Flow-CPS}(\eta>0)
  \equiv
  \text{probability-flow ODE}.
\]
这一区分直接解释了为什么 ``训练时用 CPS、评测时换 ODE'' 并不天然保持性能。

\section{为什么 RL 后的策略会依赖采样器}

\subsection{策略不是只有神经网络权重}

对固定条件 $c$，设采样器 $S$ 的第 $k$ 步转移核为
$K_{\theta,k}^{S}(x_{k-1}\mid x_k,c)$，则整条生成轨迹的分布为
\begin{equation}
  P_\theta^S(\tau\mid c)
  =p(x_T)\prod_{k=1}^{T}
    K_{\theta,k}^{S}(x_{k-1}\mid x_k,c).
  \label{eq:trajectory}
\end{equation}
理想的 sampler-conditioned 期望奖励目标可写为
\begin{equation}
  J_S(\theta)
  =\mathbb{E}_{c,\tau\sim P_\theta^S(\cdot\mid c)}
     \left[R(x_0,c)\right].
  \label{eq:sampler-objective}
\end{equation}
当前实现通过 group-relative advantage、clipped GRPO、timestep subsampling
和后述 MSE log-probability surrogate 来近似优化该目标，
并不是式 \eqref{eq:sampler-objective} 的严格无偏 policy gradient。
同一组权重 $\theta$ 搭配不同 $S$，会访问不同的中间状态分布，也会产生不同终点分布。
因此更准确的写法是两个目标泛函一般不恒等：
\[
  J_{\mathrm{CPS}}\not\equiv J_{\mathrm{ODE}},
\]
它们在某个特定 $\theta$ 上可能偶然取到相同数值，但没有一般保证。

\subsection{当前训练具体优化了什么}

历史 strict In-Domain 检查点的关键协议为：

\begin{table}[htbp]
  \centering
  \caption{历史检查点 600 的训练协议与两类评测协议}
  \label{tab:protocol}
  \small
  \begin{tabularx}{\textwidth}{>{\raggedright\arraybackslash}p{3.25cm}
                               >{\raggedright\arraybackslash}X
                               >{\raggedright\arraybackslash}X
                               >{\raggedright\arraybackslash}X}
    \toprule
    项目 & RL 训练 & CPS 评测 & 常规 ODE 评测\\
    \midrule
    转移核
      & Flow-CPS
      & Flow-CPS
      & UniPC 或 FlowMatch Euler\\
    随机系数
      & 每个 prompt/group 采样 $\eta\sim U[0,1]$
      & 固定 $\eta=0,0.3,0.7$
      & 无逐步新噪声\\
    采样步数
      & 30
      & 30
      & 50\\
    CFG
      & 1.0
      & 1.0
      & 5.0\\
    KL 系数
      & 0
      & 不适用
      & 不适用\\
    reward/scorer
      & 当时固定的旧 \code{main\_v2}
      & 最终可重评分
      & 最终可重评分\\
    \bottomrule
  \end{tabularx}
\end{table}

对应的理想化训练 mixture target 可写为
\begin{equation}
  J_{\mathrm{train}}(\theta)
  =
  \mathbb{E}_{c\sim\mathcal{C}}
  \mathbb{E}_{\eta\sim U[0,1]}
  \mathbb{E}_{\tau\sim
    P_\theta^{\mathrm{CPS}(\eta,T=30,\mathrm{CFG}=1)}
    (\cdot\mid c)}
  [R(x_0,c)].
  \label{eq:training-mixture}
\end{equation}

训练时，同一个 prompt 的 GRPO group 共用一个 $\eta$，并复用同一个初始噪声；
group 内的多样性来自 CPS 每一步重新采样的 $z$。因此网络获得的梯度是
``在随机 CPS 轨迹访问到的状态上，怎样提高 VBVR rule reward''。
固定 $\eta=0.7$ 位于训练混合分布支持集内部，$\eta=0$ 则是确定性边界极限。
连续分布中任何精确单点（包括 0.7）的概率都为零，但 0.7 周围有双侧正质量，
而 0 只在单侧邻域出现，且靠近 0 的总质量较小；精确 $\eta=0$
还会完全移除逐步探索噪声。

\begin{observation}
如果训练使用 \code{dancegrpo\_share\_group\_init\_noise=true}，
又把整个 GRPO group 强制设为严格 $\eta=0$，则同 prompt 的 rollout
会退化为相同的确定性轨迹，group advantage 缺少有效方差。
所以 ``直接把训练全部改为 $\eta=0$'' 不是可行的 GRPO 修复方案。
\end{observation}

\subsection{非线性奖励放大分布差异}

VBVR-Pro 的任务级规则包含阈值、几何关系、检测与 OCR 等非线性操作。
即使两个 sampler 的 latent 均值接近，通常也有
\[
  \mathbb{E}[R(X)]\neq R(\mathbb{E}[X]).
\]
这里并不把 ODE 输出识别为 CPS endpoint 的 $\mathbb{E}[X]$；
该式只说明逐步加入零均值噪声并不推出期望 reward 不变。
随机 CPS 可能改变轨迹覆盖范围、边缘清晰度、物体位置以及通过离散判定阈值的概率；
RL 会进一步利用这些差异。故 CPS 的收益不能简单解释成
``加噪声后取平均等于 ODE''。

\section{Wan-Trainer 实现审计}

\subsection{CPS 转移公式：一致}

核心实现位于
\path{src/models/wan_i2v.py} 的
\code{\_flowcps\_transition\_mean}。代码计算
\begin{align}
  \widehat{x}_0 &=x_t-t v_\theta,\\
  \widehat{x}_1 &=x_t+(1-t)v_\theta,\\
  \sigma_z &=s\sin(\pi\eta/2),\\
  a_{\widehat{x}_1}
    &=\sqrt{s^2-\sigma_z^2}
      =s\cos(\pi\eta/2),\\
  \mu_\theta&=(1-s)\widehat{x}_0
              +a_{\widehat{x}_1}\widehat{x}_1,
 \qquad
  x_s=\mu_\theta+\sigma_z z.
\end{align}
这与式 \eqref{eq:flow-cps} 一致。标量 $\eta$ 和 batch 级 $\eta$ 都有范围检查；
现有测试覆盖系数保持公式、每样本噪声强度、随机步和训练代理。

\subsection{CPS log-probability surrogate：代码与论文 v1 字面符号的差异}

设有效的非 batch 元素数为 $D$，项目对 CPS 使用
\begin{equation}
  \widetilde{\ell}_\theta(x_s\mid x_t)
  =-\frac{1}{D}\left\|x_s-\mu_\theta(x_t)\right\|_2^2
  =-\operatorname{mean}\!\left[(x_s-\mu_\theta(x_t))^2\right],
  \label{eq:surrogate-logp}
\end{equation}
对应代码
\code{\_flowcps\_transition\_log\_prob}。
当 $\sigma_z>0$ 时，各向同性 $D$ 维 Gaussian 的严格 log-density 应为
\begin{equation}
  \log p_\theta(x_s\mid x_t)
  =-\frac{\|x_s-\mu_\theta\|_2^2}{2\sigma_z^2}
   -D\log\sigma_z-\frac{D}{2}\log(2\pi).
  \label{eq:strict-gaussian-logp}
\end{equation}
因此必须精确区分：

\begin{itemize}
  \item arXiv v1 的相邻印刷公式把平方误差写成正号；若它被称为
  log-probability，标准概率意义要求负号。当前代码与合理的负误差解释一致，
  但与 v1 公式的字面符号不一致，最可能是论文排版错误。
  \item 同一 timestep 的 old/current ratio 中，若 $\sigma_z$ 固定，
  式 \eqref{eq:strict-gaussian-logp} 的加性归一化项可以相消；
  $1/(2\sigma_z^2)$ 是乘法尺度，不能用 ``相消'' 解释其省略。
  省略该因子会改变不同 $\eta$ 与 timestep 之间的相对加权。
  \item 代码还使用 $1/D$ 的 mean reduction：在固定 latent shape 下，
  它是共同的全局尺度，会重标定 PPO ratio/clip 行为；
  跨不同 latent 维数时，它使尺度按维数归一化。
  \item 当 $\sigma_z=0$ 时（包括 $\eta=0$ 的所有步骤，以及任意
  $\eta$ 的最终 $s=0$ 步），转移退化为 Dirac measure，
  不存在相对于普通 Lebesgue measure 的有限 Gaussian density；
  此时代码返回的有限 MSE 只能视作 heuristic surrogate。
\end{itemize}

所以这里发现的是\textbf{论文 v1 字面公式与代码的符号不一致}，
不是 CPS transition 公式抄错。该 surrogate 是否促进了 sampler co-adaptation
或分数 gap，仍需用归一化尺度和 $\eta$ 分层消融来判断。

\subsection{每组 \texorpdfstring{$\eta$}{eta} 的采样与 replay：闭环}

\path{src/trainer/dancegrpo_trainer.py} 的
\path{_sample_group_cps_noise_levels} 对每个 prompt/group
生成一个可复现的 $\eta$；同组 $G$ 条 rollout 共用该值。
rollout 保存 \code{cps\_noise\_levels}，训练 replay 再把同一系数传回
transition mean/log-probability。这样避免了 rollout 用一个 $\eta$、
policy ratio 重算却用另一个 $\eta$ 的严重不一致。

\subsection{真正存在的评测不配平}

常规 ODE wrapper 固定 $T=50$、CFG $=5$；CPS wrapper 固定
$T=30$、CFG $=1$，并调用训练时 \code{flowcps} 路径。
因此原始比较一次改变了至少三个变量：

\begin{equation}
  \underbrace{\text{sampler}}_{\mathrm{CPS}\leftrightarrow\mathrm{ODE}}
  +\underbrace{\text{time grid}}_{30\leftrightarrow50}
  +\underbrace{\text{CFG}}_{1\leftrightarrow5}.
\end{equation}

其中 CFG 尤其值得注意。RL 训练只在 CFG $=1$ 的有效速度场上优化；
评测时把 CFG 提到 5，会改变每一步 conditional/unconditional
组合后的速度，可能造成过度引导、轨迹偏移或规则任务中的几何失真。
所以不能把 UniPC/ODE 的全部劣势归因于 ``ODE 理论上更差''。

\subsection{视频与 scorer 输入链路：已经对齐}

训练 reward 与最终生成评测共享以下契约：

\begin{itemize}
  \item 保留 161 帧，不裁成 81 帧；
  \item 空间缩放到 $1024\times1024$，不做内容裁剪；
  \item 通过提高编码 FPS 把播放时长限制在 5 秒内，161 帧使用至少 33 FPS；
  \item 使用相同的 Diffusers video export 语义；
  \item 把 prompt、GT path 和 metadata 交给同一个
  \code{run\_evaluation.evaluate\_single\_video} 入口；
  \item scorer 运行在隐藏 CUDA 的 CPU 子进程中，避免 EasyOCR 占用训练 GPU。
\end{itemize}

G-15 与 EasyOCR-heavy 的 O-34 已做训练路径/最终路径双路对齐验证：
raw 与 prepared 帧一致，隔离 scorer 的结果一致。该验证降低了 I/O 差异风险，
但只覆盖两个代表任务，不能替代 100 个任务的完整语义审计。

\section{EvalKit 更新审计}

\subsection{版本变化}

历史评测固定在 GitHub \code{main\_v2} revision
\code{42a1593d8e493370c768be8e43646f0e0a9d8525}。
截至 2026-07-24，\textbf{可验证的 GitHub \code{main\_v2} latest}
\cite{evalkit} 为：
\begin{center}
  \small\code{6fedd9d9edb8daafa56aca8e53885aa8ad6f6037}
\end{center}
它与旧版相隔 14 个提交。比较范围包含 13 个文件，共
1,756 行新增、359 行删除；主要变化位于十个 evaluator 分片、
\code{evaluators/\_\_init\_\_.py}、base evaluator 和 utilities。
runner 与 requirements 没有变化，注册任务仍为 100 个。
GitHub 默认 \code{main} 是另一条历史线，不能用其 HEAD 替代
\code{main\_v2}；受访问权限限制，本文也没有独立核验 GitLab 上游状态。

完整提交比较见
\href{https://github.com/xujunxiangwork/VBVR-Evalkit-Interleave/compare/42a1593d8e493370c768be8e43646f0e0a9d8525...6fedd9d9edb8daafa56aca8e53885aa8ad6f6037}
{GitHub compare view}。

\subsection{语义变化为何会改分}

更新不是纯重构，而包含实际评分语义修复，例如：

\begin{itemize}
  \item G-15 对移动障碍物行为新增/强化惩罚；
  \item 多个任务把 consistency 从加分项改成 penalty multiplier；
  \item 修复 object 与 annotation 混淆、方向性 resize、黑色 RGB 阈值；
  \item 增加 convex-hull fallback、圆/环检测、HSV 与 ray-length 修复；
  \item 对空白输出移除基础分，调整 shape cliffs 和若干任务阈值；
  \item 修复 G-194、G-140、O-9 至 O-14、O-18/O-19 等具体 evaluator。
\end{itemize}

这些变化意味着相同视频在新旧 scorer 下可以得到不同分数。
因此所有报表必须把 \textbf{EvalKit revision + 完整 scorer source hash}
作为结果合同的一部分。

\subsection{新的合同 hash 与代码更新}

最新固定版本使用：
\begin{center}
  revision\\[-0.2em]
  \small\code{6fedd9d9edb8daafa56aca8e53885aa8ad6f6037}\\[0.4em]
  source SHA-256\\[-0.2em]
  \small\code{eb977da60e95456734063ba018b14d805680179fdf0e3e3b2ba6f603f27a935c}
\end{center}
后者覆盖 runner、所有 VBVR Python 源文件、75 个 annotation 文件和
requirements，比只记录 Git commit
更能防止 checkout 后被本地修改而不自知。

本次代码更新形成以下闭环：

\begin{enumerate}
  \item 训练配置必须提供 64 位十六进制 scorer hash；
  \item reward 初始化时计算实际 hash，不匹配立即失败；
  \item 最终评分结果记录 revision、source hash、输入/输出 manifest；
  \item sweep 在跳过已有结果前验证 500 视频、500 无错误样本、100 任务、
  完整 provenance 和 scorer hash；
  \item 汇总工具发现多个 scorer fingerprint 时直接拒绝，
  阻止新旧分数静默混入同一趋势表；
  \item latest scorer 使用独立 checkout、checkpoint/output/W\&B/tmp 命名空间。
\end{enumerate}

\warnbox{\textbf{关键历史事实：}
检查点 600 是用旧 scorer reward 训练出来的；本报告的 ``latest''
数字是以各模式已有的 prepared videos 为输入、使用新 scorer 的复评分。
这能回答 ``当前 EvalKit 如何看这些视频''，但不能声称
``模型已经使用最新 reward 训练''。真正的 latest-reward 训练应使用新增的
\path{configs/train_dancegrpo_vbvr_pro_5b_256x256x161_rule_cps_from_nsft_bs_32_lr_1e-6_manifest_rl_evalkit_6fedd9d9.yaml}。}

\section{检查点 600：最新 EvalKit 实验结果}

\subsection{评测完整性}

四种模式均满足：

\begin{itemize}
  \item 500 个生成/准备后视频；
  \item 500 个成功评分样本，错误数为 0；
  \item 100 个任务；
  \item 同一 latest EvalKit revision 与 source hash；
  \item 同一 split manifest 命名空间；
  \item score provenance 为 complete，scorer identity、结果路径绑定、
  结果 JSON 当前指纹与样本/任务计数校验通过。
\end{itemize}

这里没有声称重算了 provenance 的全部输入树：
本文分析脚本只重新验证 \code{output\_files} 中的结果文件。
CPS 两个模式的旧 scorer wrapper 在评分后已变化，因而完整 all-section
provenance 不能用当前 checkout 重新通过；这不影响已固定结果 JSON 的读取，
但限制了 ``全输入可重放'' 的证据等级。

\subsection{四种采样模式}

\begin{table}[htbp]
  \centering
  \caption{检查点 600 在最新 EvalKit 上的结果。粗绿色为每列最优。}
  \label{tab:latest-scores}
  \small
  \begin{tabular}{lcccccc}
    \toprule
    采样模式 & $T$ & CFG & $\eta$ & Overall & In-Domain & Out-of-Domain\\
    \midrule
    UniPC ODE             & 50 & 5.0 & --  & 0.504266 & 0.628499 & 0.380033\\
    FlowMatch Euler ODE   & 50 & 5.0 & --  & 0.504909 & 0.634534 & 0.375284\\
    Flow-CPS              & 30 & 1.0 & 0.0 & 0.510376 & 0.630510 & 0.390241\\
    Flow-CPS              & 30 & 1.0 & 0.7 & \best{0.534130} & \best{0.666700} & \best{0.401560}\\
    \bottomrule
  \end{tabular}
\end{table}

从表 \ref{tab:latest-scores} 可以得到三层结论：

\begin{enumerate}
  \item \textbf{换高阶 ODE solver 没有消除差距。}
  UniPC 与一阶 Euler 的 Overall 只差 $-0.000644$；
  至少在当前 $T=50$、CFG $=5$ 设置下，数值阶数不是主要瓶颈。
  \item \textbf{评测协议确实影响结果。}
  $\eta=0$、$T=30$、CFG $=1$ 的控制比 Euler ODE
  高 $0.005466$，但该差值同时包含 pipeline/time-grid/CFG 变化，
  不能只归因于 CPS API。
  \item \textbf{存在统计稳定的 sampler-associated signal。}
  在同一个 CPS 实现、同样 $T=30$ 和 CFG $=1$ 下，
  $\eta=0.7$ 比 $\eta=0$ 高 $0.023754$。
  这是当前最接近 ``只改变 stochasticity'' 的观察，但由于两套视频没有
  initial-latent content hash，仍不是严格单变量因果实验。
\end{enumerate}

\subsection{任务级 task-paired bootstrap}

对 100 个任务的平均分差进行有放回重采样，共 100,000 次，
随机种子为 20260724，报告 percentile 95\% 区间。

\begin{table}[htbp]
  \centering
  \caption{CPS $\eta=0.7$ 相对控制方法的任务级差值}
  \label{tab:bootstrap}
  \small
  \begin{tabular}{lcccc}
    \toprule
    控制方法 & Overall 差值 & 任务均值差 95\% CI & 胜/平/负任务 & 中位任务差\\
    \midrule
    UniPC ODE
      & +0.029864 & [0.006433,\ 0.053919] & 60/7/33 & +0.009341\\
    FlowMatch Euler ODE
      & +0.029220 & [0.007166,\ 0.051996] & 61/9/30 & +0.008072\\
    Flow-CPS $\eta=0$
      & +0.023754 & [0.003952,\ 0.044239] & 58/9/33 & +0.001062\\
    \bottomrule
  \end{tabular}
\end{table}

三个区间都不跨零，说明在任务级重采样下平均差稳定为正。
但这并不排除少数大增益任务主导均值；中位任务差明显小于均值差，
说明收益分布存在长尾：
部分任务获得较大提升，另一些任务退化。下一步应结合
\path{scripts/dev/collect_vbvr_task_regressions.py}
查看哪些 evaluator/任务承担了主要增益与损失。

\warnbox{\textbf{统计解释边界：}
这里的 paired 单位是 ``相同任务名称''，不是 ``相同初始 latent 的成对视频''。
相同 seed 文本也不足以证明两个不同 pipeline 生成了逐元素相同的 $x_T$。
因此置信区间说明任务均值差稳定，但不能独自建立 sampler 的严格因果效应。}

\subsection{新旧 scorer 迁移量}

\begin{table}[htbp]
  \centering
  \caption{同一批模式从旧 scorer 到最新 scorer 的分数变化}
  \label{tab:migration}
  \small
  \begin{tabular}{lrrrrrr}
    \toprule
    & \multicolumn{2}{c}{Overall}
    & \multicolumn{2}{c}{In-Domain}
    & \multicolumn{2}{c}{Out-of-Domain}\\
    \cmidrule(lr){2-3}\cmidrule(lr){4-5}\cmidrule(lr){6-7}
    模式 & 旧 & 新 & 旧 & 新 & 旧 & 新\\
    \midrule
    UniPC ODE
      & 0.508712 & 0.504266
      & 0.612457 & 0.628499
      & 0.404967 & 0.380033\\
    Euler ODE
      & 0.510049 & 0.504909
      & 0.619022 & 0.634534
      & 0.401075 & 0.375284\\
    CPS $\eta=0$
      & 0.512337 & 0.510376
      & 0.627383 & 0.630510
      & 0.397290 & 0.390241\\
    CPS $\eta=0.7$
      & 0.534556 & 0.534130
      & 0.656026 & 0.666700
      & 0.413086 & 0.401560\\
    \bottomrule
  \end{tabular}
\end{table}

迁移方向很一致：最新 scorer 在这四种模式上都提高了 In-Domain，
同时降低了 Out-of-Domain，且 Overall 四项均略降。
UniPC、Euler、CPS $\eta=0$、CPS $\eta=0.7$ 的 Overall 新减旧分别为
$-0.004446$、$-0.005139$、$-0.001961$、$-0.000426$；
对应 ID 变化为 $+0.016042$、$+0.015512$、$+0.003127$、$+0.010674$，
OOD 变化为 $-0.024934$、$-0.025791$、$-0.007049$、$-0.011526$。
这进一步说明，不能把旧 scorer 的历史曲线与 latest scorer 的新点直接连线。

旧/新 provenance 中，UniPC、Euler 和 CPS $\eta=0.7$ 的
prepared-video tree fingerprint 不相同；条目数和总字节数一致，但当前
\code{fingerprint\_tree} 把 mtime 纳入摘要，差异可能仅来自时间戳，
也不能反向证明 500 个 MP4 字节完全相同。CPS $\eta=0$ 的 tree fingerprint
一致。因此本表是\textbf{相同模式/路径命名下的 scorer migration}，
而不是四组都具备 content-only SHA 证明的严格逐视频配对迁移。

\section{对 ``为什么 CPS 明显优于 ODE'' 的综合解释}

\subsection{候选机制与当前证据强弱}

结合公式、代码和实验，当前可按证据强弱排列候选机制，但不能把这一排序
当作已经完成的因果分解：

\begin{enumerate}
  \item \textbf{sampling-protocol dependence（gap 证据强，RL 因果证据中等）。}
  策略在随机 CPS 轨迹的状态分布上接受 reward/advantage 更新，
  固定 $\eta=0.7$ 位于训练 mixture 的内部支持；custom $\eta=0$ 是
  同一 CPS kernel family 的确定性边界极限，并非完全无关的 kernel。
  明显未按训练协议采样的是 $T=50$/CFG 5 的 Diffusers ODE。
  当前 post-RL gap 与 sampler co-adaptation 一致，但也可能部分继承自
  base/SFT 模型；需比较 pre-RL 与 post-RL 的 sampler gap，
  做 difference-in-differences，最好再覆盖多 seed/checkpoint。
  \item \textbf{CFG mismatch（高可信、尚待单变量实验证实）。}
  训练/CPS eval 使用 CFG 1，ODE eval 使用 CFG 5。
  RL 后的速度场不保证在强 CFG 外推下仍保持规则任务结构。
  \item \textbf{时间网格与 pipeline mismatch（中高可信）。}
  $T=30$ 的内部 CPS 路径与 $T=50$ 的 Diffusers ODE 路径不完全可比。
  \item \textbf{该检查点在随机 CPS 下更易通过 reward 阈值（中可信）。}
  $\eta=0.7$ 相对同路径 $\eta=0$ 的 $+0.023754$ 支持这一点；
  可能来自轨迹探索、局部几何或阈值通过率的改变，但尚无同 latent
  单变量生成来确定 stochasticity 的因果贡献。
  \item \textbf{ODE solver 阶数（当前证据较弱）。}
  UniPC 与 Euler 分数几乎相同，表明仅换高阶 solver 不足以解决。
\end{enumerate}

\subsection{哪些说法目前不成立}

\begin{itemize}
  \item 不能说 ``CPS 理论上应与 ODE 同分''；论文没有该定理。
  \item 不能说 ``$\eta=0$ 结果证明 CPS 和所有 ODE pipeline 等价''；
  等价只在相同网格、速度场和实现细节下成立。
  \item 不能说 ``最新 EvalKit 让 CPS 优势出现''；
  旧 scorer 下 CPS $\eta=0.7$ 同样领先，更新只改变了绝对分数和任务权重。
  \item 不能说 ``Out-of-Domain 是 zero-shot''；任务类型在 RL 数据中出现过。
  \item 不能把同一整数 seed 当作同一初始 latent 的证明。
\end{itemize}

\section{怎样把 ODE 分数做高}

\subsection{第一阶段：不训练，先做严格的 ODE 推理搜索}

最便宜且信息量最高的是一个\textbf{配平的生成矩阵}。所有 arm 必须保存相同
initial latent tensor，记录其 SHA-256，并统一 condition、VAE、dtype、
flow shift、输出视频与 scorer：

\infobox{\textbf{这是下一步实验设计，不是当前已经可直接运行的完整命令。}
现有 CPS CLI 在内部按 seed 采样 latent，Diffusers CLI 只接收 generator，
二者都没有跨 pipeline 的 latent export/injection 参数；也没有 strict batch
的内部 Euler wrapper 与统一 timestep dump。实施表 \ref{tab:matched-grid}
前，至少要补 latent 导出/注入及 SHA-256、逐步 latent dump，
以及批量内部 Euler 入口。}

\begin{table}[htbp]
  \centering
  \caption{建议的第一轮配平实验}
  \label{tab:matched-grid}
  \small
  \begin{tabularx}{\textwidth}{clcccX}
    \toprule
    优先级 & 采样器 & $T$ & CFG & 随机性 & 目的\\
    \midrule
    P0 & Flow-CPS $\eta=0.7$ & 30 & 1 & 有
       & 训练匹配基准\\
    P0 & Flow-CPS $\eta=0$ & 30 & 1 & 无
       & 已有近似 deterministic 控制\\
    P0 & 内部显式 Euler & 30 & 1 & 无
       & 与 $\eta=0$ 检查逐步误差是否在容差内及其累积\\
    P0 & Diffusers Euler & 30 & 1 & 无
       & 分离 pipeline 差异\\
    P1 & UniPC & 30 & 1 & 无
       & 只改变 solver 阶数\\
    P1 & Euler/UniPC & 50 & 1 & 无
       & 只改变步数\\
    P1 & Euler/UniPC & 30 & 2,3,5 & 无
       & 定位 CFG 过强问题\\
    P2 & Euler/UniPC & 20,40,50,75 & 最优 CFG & 无
       & 在验证集上搜索时间网格\\
    \bottomrule
  \end{tabularx}
\end{table}

推荐判据：

\begin{enumerate}
  \item 先在不使用 bench 选择超参数的验证样本上调 $T$/CFG；
  \item 固定一个方案后，只在 500 样本 bench 上做一次确认；
  \item 同时报 Overall、ID、OOD、100-task paired delta 和视频失败率；
  \item 除平均分外，检查 CPS 优势最大的任务是否由 stochasticity
  改善，还是由 ODE CFG 导致几何/颜色退化。
\end{enumerate}

从现有数据判断，\textbf{把 ODE 的 CFG 从 5 降到接近训练值 1}
是最值得先试的单项改动。单纯从 Euler 换 UniPC 已经没有显示出收益。

\subsection{第二阶段：加入 ODE-aware 训练，而不牺牲 GRPO 探索}

若配平后 ODE 仍明显落后，可把训练目标改为混合策略：
\begin{equation}
  J_{\mathrm{mix}}(\theta)
  =\lambda J_{\mathrm{CPS}}(\theta)
   +(1-\lambda)J_{\mathrm{ODE}}(\theta)
   -\beta\,\mathcal{L}_{\mathrm{reference\text{-}mean}}.
  \label{eq:mixed-objective}
\end{equation}
可实施的路线按风险从低到高为：

\begin{enumerate}
  \item \textbf{ODE endpoint monitor。}
  每隔若干 step 在固定小验证集上做 $\eta=0$/Euler 生成，只监控不反传，
  用于早停和选 checkpoint。
  \item \textbf{reference-mean consistency/MSE surrogate 消融。}
  当前 $\beta=0$，可能允许模型为 CPS reward 过拟合并损伤原基座的 ODE 场。
  当前 Flow-CPS 代码所谓 ``KL'' 实际是 policy/reference transition mean
  的未归一化 MSE，没有除以 CPS 方差，并非严格 Gaussian KL。
  应先记录 reward、policy loss 与该 MSE 的尺度，再做对数级 $\beta$ 搜索；
  $\{10^{-4},5\times10^{-4},10^{-3},4\times10^{-3}\}$ 只能作为待校准起点，
  不能直接按论文 KL 系数解释。论文在部分可验证奖励上也观察到无 KL 会退化。
  \item \textbf{CPS--ODE consistency auxiliary loss。}
  对相同 $x_t$，约束 $\eta=0$ 的一步/终点与 teacher ODE 或参考模型接近，
  同时保留 $\eta>0$ rollout 提供 group 多样性。
  \item \textbf{双 sampler reward。}
  对同一 prompt 同时评 CPS 和 deterministic endpoint，
  优化二者加权均值或较差者，避免只提升一条采样路径。
  \item \textbf{CFG-aware rollout。}
  如果生产必须用 CFG $>1$，训练时随机化 CFG 或加入相应路径；
  否则评测应优先采用训练匹配 CFG。
\end{enumerate}

上述 consistency、双 sampler reward 与 CFG-aware rollout
目前都是设计建议，不是已存在的配置开关。尤其随机 CFG 需要像
\code{cps\_noise\_levels} 一样把每组 CFG 写入 rollout，并在 replay 中恢复，
不能只在配置中改一个数值。

\warnbox{\textbf{不建议：}
不要把所有 GRPO rollout 直接改为 $\eta=0$ 并继续共用 group initial noise。
这会让同 prompt 的多条轨迹完全相同，advantage 标准化退化。
若要让确定性 ODE 直接参与优化，需要不同初始噪声、离策略 endpoint loss、
蒸馏/一致性目标，或与 $\eta>0$ 的探索 rollout 混合。}

\subsection{第三阶段：用 latest scorer 重新训练}

新训练配置已经将以下变量固定并隔离：

\begin{itemize}
  \item manifest RL：50 个 In-Domain 任务、50,000 个样本；
  \item 固定 Flow-CPS $\eta=0.7$、$T=30$、CFG $=1$；
  \item EvalKit revision \code{6fedd9d9...} 与完整 source hash；
  \item 独立 checkpoint、W\&B、reward tmp 和 dataset descriptor；
  \item 从现有 BF16 Diffusers 转换模型初始化，\code{resume\_from: null}。
\end{itemize}

最后一点意味着它是\textbf{权重初始化}而不是历史 DCP optimizer/dataloader
状态续训；并且无法恢复此前 BF16 转换已经丢失的精度。
必须更正验证边界：已落盘的 8xH100 manifest-RL 一步 smoke
使用的是旧 scorer hash \code{f99edd06...}，不是 latest
\code{eb977da6...}。latest 配置目前只完成了解析、50,000 样本数据集和
静态/单元链路验证；其 checkpoint 与 reward tmp 目录尚不存在，
没有 optimizer step，也没有启动长训。

此外，latest 配置写了：
\begin{center}
  \code{grpo\_sample\_batch\_size: 8}\\
  \code{grpo\_train\_sample\_batch\_size: 4}
\end{center}
但当前 shared-prompt
实现按 rollout batch 8 建 chunk，replay 直接复用相同 chunk；
所以实际 replay micro-batch 是 8 而不是 4。这是 latest 长训前必须修复
或以真实显存 smoke 验证的风险，不能把生产配置视为已经闭环。

\section{推荐的结论判定表}

\begin{longtable}{p{4.1cm}p{5.4cm}p{5.4cm}}
  \caption{后续实验结果应如何解释}
  \label{tab:decision}\\
  \toprule
  观察 & 更可能的原因 & 下一步\\
  \midrule
  \endfirsthead
  \toprule
  观察 & 更可能的原因 & 下一步\\
  \midrule
  \endhead
  Euler $T=30$/CFG 1 与 CPS $\eta=0$ 的逐步误差超出预设容差或持续累积
    & pipeline、time grid、scheduler 或实现细节不一致
    & 做逐步 latent diff；先修数值合同，不讨论 reward gap\\
  \addlinespace
  二者在容差内一致，但 UniPC/CFG 5 明显差
    & solver/CFG 改变了 RL 后速度场的访问分布
    & 以 CFG 1 为中心调参；分开报告 solver 与 CFG\\
  \addlinespace
  所有 deterministic 设置都比 CPS $\eta=0.7$ 差
    & sampler-associated gap；可能是 base gap、RL co-adaptation 或二者共同作用
    & 比较 pre/post-RL gap；再上 ODE consistency、双 sampler objective、reference MSE\\
  \addlinespace
  latest scorer 下绝对分变，采样器排序稳定
    & scorer 语义改变但 sampler gap 稳健
    & 保留版本隔离；不可拼接新旧曲线\\
  \addlinespace
  latest-reward 重训后 ODE gap 缩小
    & 旧 reward 过拟合或新 evaluator 更平滑
    & 扩展到多 checkpoint、多 seed 验证\\
  \addlinespace
  latest-reward 重训后 CPS gap仍在
    & sampler-conditioned policy 的解释更强，但仍需 pre/post 与多 seed
    & 根据部署 sampler 决定是否接受，或显式多 sampler 训练\\
  \bottomrule
\end{longtable}

\section{复现入口}

\subsection{重算本文数字}

本文的结果 JSON 验证、旧/新 scorer 迁移量和任务级 bootstrap
由以下脚本重算：

\begin{verbatim}
.venv/bin/python scripts/dev/analyze_cps_ode_checkpoint600.py
\end{verbatim}

默认输出：
\begin{center}
  \path{storage/reports/cps_ode_evalkit_report/checkpoint600_metrics.json}
\end{center}

脚本会验证每个 latest provenance 的 complete 状态、scorer identity、
结果路径绑定、结果 JSON 当前指纹、500 个无错误样本、100 个任务、
revision 与 source hash；任一条件不满足都会失败。它有意只重算
\code{output\_files}，不宣称重新验证 generation/preparation/scorer-wrapper
等全部输入树。

\subsection{主要源文件}

\begin{description}
  \item[CPS 数学与 rollout]
  \path{src/models/wan_i2v.py}，
  \path{src/trainer/dancegrpo_trainer.py}
  \item[CPS 评测入口]
  \path{src/cli/eval_i2v_cps.py}
  \item[reward 对齐]
  \path{src/trainer/rewards/vbvr_rule.py}
  \item[scorer 与结果合同]
  \path{src/eval/vbvr_run_evaluation_parallel.py}\\
  \path{src/eval/evaluation_provenance.py}
  \item[主评测 launcher]
  \path{scripts/eval/vbvr_pro/vbvr_pro_5b_main_v2.fish}
  \item[strict sweep]
  {\footnotesize\path{scripts/eval/vbvr_pro/dancegrpo_indomain_strict/vbvr_pro_5b_dancegrpo_indomain_strict_sweep_main_v2.fish}}
  \item[latest 训练配置]
  \path{configs/train_dancegrpo_vbvr_pro_5b_256x256x161_rule_cps_from_nsft_bs_32_lr_1e-6_manifest_rl_evalkit_6fedd9d9.yaml}
\end{description}

\subsection{验证状态}

一次完整项目测试执行返回 \textbf{130 passed, 1 warning}，但没有单独保存
该次 pytest 全量日志；\code{.pytest\_cache} 不能替代运行 provenance。
Python lint 与变更过的 fish launcher 语法检查通过。
EvalKit 安装锁 smoke 验证了 owner/follower 等待、原子发布与清锁，
但它故意传入错误期望 hash 后终止，不代表一次完整评测成功。
G-15 和 O-34 的 validator 比较的是视频解码后的 RGB frame arrays，
并报告 raw/prepared 帧逐元素一致；它没有比较 MP4 文件字节或 content SHA。
当次 validator 也报告 scorer 分数一致，但分数 stdout/JSON 没有独立保存。

\section{局限性}

\begin{enumerate}
  \item 当前 latest scorer 结果只覆盖检查点 600，尚未重评分完整 checkpoint sweep。
  \item 复评分使用了已存在的 prepared videos，不是从模型生成开始的全新端到端
  latest run；3/4 模式的新旧 tree fingerprint 不同且含 mtime，
  尚无全量 MP4 content-only hash 来证明严格逐视频不变。
  \item UniPC/Euler 使用 $T=50$/CFG 5，CPS 使用 $T=30$/CFG 1；
  只有 CPS $\eta=0$ 与 $\eta=0.7$ 是较接近的控制。
  \item 跨采样器没有保存并核对同一 initial latent tensor 的 hash。
  \item bootstrap 以 100 个任务为单位；每个任务只有 5 个 bench 样本，
  对任务内方差的估计有限。
  \item checkpoint 600 是旧 scorer reward 训练产物；
  latest scorer 复评分不能替代 latest-reward 重训。
  \item 没有 pre-RL/SFT 同协议 sampler gap，因而不能对 RL 造成的增量做
  difference-in-differences；目前只有 post-RL checkpoint。
  \item latest-reward 配置尚无 optimizer-step smoke 或长训，
  且当前 shared-prompt replay 实际沿用 rollout chunk 8，
  配置的 replay micro-batch 4 尚未生效。
  \item reward I/O 对齐验证只抽查 G-15 和 O-34。
  \item 训练数据包含全部 100 个受支持任务类型，所谓 OOD 不是零样本任务泛化。
\end{enumerate}

\section{最终结论}

本次审计没有发现 Wan-Trainer 的 Flow-CPS \textbf{转移公式}
与论文之间的实质性不一致；但 arXiv v1 的 log-probability 公式字面正号
与代码的负平方误差不一致，按概率意义最可能是论文排版错误，
而代码本身采用的仍是未归一化 heuristic surrogate。
当所有条件完全相同时，$\eta=0$ 的 Flow-CPS 在代数上就是一阶 Euler；
但 $\eta>0$ 的 CPS 与 ODE 并不共享逐路径或有限步终点等价性。

检查点 600 的数据同时支持两个判断：

\begin{enumerate}
  \item 原始 CPS--ODE 比较被 $T=30$ vs.\ $50$、CFG 1 vs.\ 5 和不同 pipeline
  混杂，ODE 分数存在通过匹配训练 CFG、搜索时间网格和严格统一 latent 而提升的空间；
  \item 即便在同一 CPS 路径内配平 $T$/CFG，$\eta=0.7$ 仍比 $\eta=0$
  高 $0.023754$，任务级 95\% bootstrap 区间为
  $[0.003952,0.044239]$，说明当前样本集存在统计稳定的
  sampler-associated gap；它与 RL co-adaptation 一致，但还不能排除
  base/SFT 模型原有 gap。
\end{enumerate}

因此，短期最优行动不是简单地 ``把 ODE 换成更高阶 solver''，
而是先做同 latent、同 CFG、同时间网格的矩阵，优先测试 CFG 1 的 ODE；
若差距仍在，再通过 reference-mean MSE、CPS--ODE consistency 或双 sampler reward
把 ODE 端点显式纳入训练目标。对于只追求当前 VBVR-Pro 得分的部署，
CPS $\eta=0.7$ 是\textbf{检查点 600、当前四个已审计 arm、这 500 个 bench
样本和 latest scorer 下}观测到的最佳模式；对于必须使用确定性 ODE 的部署，
则需要把 ODE 当成目标策略来训练和选择 checkpoint，而不能假设 CPS RL
会自动迁移过去。

\appendix

\section{关键数值的机器可读摘要}

\begin{verbatim}
checkpoint: 600
evalkit_revision:
  6fedd9d9edb8daafa56aca8e53885aa8ad6f6037
evalkit_source_sha256:
  eb977da60e95456734063ba018b14d805680179fdf0e3e3b2ba6f603f27a935c

latest overall:
  unipc_ode       0.5042658667104944
  euler_ode       0.5049093773490485
  flowcps_eta0    0.5103755288562246
  flowcps_eta0.7  0.5341297677135947

bootstrap:
  unit       task
  tasks      100
  resamples  100000
  seed       20260724
\end{verbatim}

\section{论文与实现映射速查}

\begin{table}[htbp]
  \centering
  \small
  \begin{tabularx}{\textwidth}{p{3.3cm}p{5.1cm}X}
    \toprule
    概念 & 数学对象 & Wan-Trainer 入口\\
    \midrule
    预测干净 latent
      & $\widehat{x}_0=x_t-t v_\theta$
      & \code{\_flowcps\_transition\_mean}\\
    预测噪声 latent
      & $\widehat{x}_1=x_t+(1-t)v_\theta$
      & \code{\_flowcps\_transition\_mean}\\
    新噪声标准差
      & $s\sin(\pi\eta/2)$
      & \code{std\_dev\_t}\\
    预测噪声系数
      & $s\cos(\pi\eta/2)$
      & \code{sqrt(sigma\_prev**2-std**2)}\\
    CPS 代理 log-prob
      & $-D^{-1}\|x_s-\mu_\theta\|^2$
      & \code{\_flowcps\_transition\_log\_prob}\\
    group 级 $\eta$
      & 每 prompt 一个值，组内共享
      & \code{\_sample\_group\_cps\_noise\_levels}\\
    final rule reward
      & prepared video $\rightarrow$ task evaluator
      & \code{VBVRRuleReward}\\
    scorer 合同
      & revision + full source SHA-256
      & \code{evaluation\_provenance.py}\\
    \bottomrule
  \end{tabularx}
\end{table}

\begin{thebibliography}{9}

\bibitem{flowcps}
Feng Wang and Zihao Yu.
\newblock \emph{Coefficients-Preserving Sampling for Reinforcement Learning with Flow Matching}.
\newblock arXiv:2509.05952v1, 2025.
\newblock \url{https://arxiv.org/html/2509.05952v1}.

\bibitem{rectifiedflow}
Xingchao Liu, Chengyue Gong, and Qiang Liu.
\newblock \emph{Flow Straight and Fast: Learning to Generate and Transfer Data with Rectified Flow}.
\newblock arXiv:2209.03003, 2022.

\bibitem{ddim}
Jiaming Song, Chenlin Meng, and Stefano Ermon.
\newblock \emph{Denoising Diffusion Implicit Models}.
\newblock ICLR, 2021.

\bibitem{unipc}
Wenliang Zhao, Lujia Bai, Yongming Rao, Jie Zhou, and Jiwen Lu.
\newblock \emph{UniPC: A Unified Predictor-Corrector Framework for Fast Sampling of Diffusion Models}.
\newblock NeurIPS, 2023.

\bibitem{evalkit}
xujunxiangwork.
\newblock \emph{VBVR-Evalkit-Interleave}, branch \code{main\_v2}.
\newblock Latest audited revision
\code{6fedd9d9edb8daafa56aca8e53885aa8ad6f6037}.
\newblock \url{https://github.com/xujunxiangwork/VBVR-Evalkit-Interleave}.

\end{thebibliography}

\end{document}
